% Generated draft: requires mathematical and source review.
\documentclass[11pt]{article}
\usepackage[margin=25mm]{geometry}
\usepackage{fontspec}
\setmainfont{Latin Modern Roman}
\usepackage{amsmath,amssymb,mathtools,unicode-math}
\setmathfont{Latin Modern Math}
\usepackage{enumitem,needspace,xurl,hyperref}
\setlength{\parindent}{0pt}
\setlength{\parskip}{5pt}
\setlength{\emergencystretch}{3em}
\newcommand{\N}{\mathbb{N}}
\newcommand{\Z}{\mathbb{Z}}
\newcommand{\Q}{\mathbb{Q}}
\newcommand{\R}{\mathbb{R}}
\newcommand{\C}{\mathbb{C}}
\newcommand{\F}{\mathbb{F}}
\newcommand{\A}{\mathbb{A}}
\newcommand{\PP}{\mathbb{P}}
\newcommand{\E}{\mathbb{E}}
\newcommand{\eps}{\varepsilon}
\newcommand{\dd}{\,\mathrm{d}}
\begin{document}
\section*{UnsolvedMath ID: 20002292}
\textbf{Source problem code:} catalogue entry
\subsection{Definitions and mathematical statement}
Pasten's Question 8 asks: is an explicitly known recursively enumerable subset of \(\Q\) non-Diophantine over \(\Q\)? What are the thinnest known Diophantine rational sets? Does a suitable Bombieri--Lang hypothesis force such sets to be finite or to have fast growth, potentially proving \(\Z\) non-Diophantine? The source gives no growth parameter or precise ``fast'' threshold. A Diophantine set means an existential polynomial projection over \(\Q\); recursively enumerable (listable) uses a computable encoding of rational numbers, with no decidability of membership required.

The aid reports that existence of a listable non-Diophantine subset of \(\Q\) is equivalent to \(\Z\) not being Diophantine in \(\Q\), and remains open in its supplied assessment. It also reports no known general finite-or-fast height-growth dichotomy. The original statement of evidence against a Diophantine definition of \(\Z\) is historical context, not a proof.

For its unconditional height calibration, write a reduced rational as \(x=a/b\), \(b>0\), and set
\[
 H(x)=\max(|a|,b),\qquad N_S(B)=\#\{x\in S:H(x)\le B\}.
\]
Consider
\[
 E:y^2+y=x^3-x,\qquad P=(0,0),\qquad
 S=\{x\in\Q:\exists y\in\Q,\ y^2+y=x^3-x\}.
\]
This \(S\) is Diophantine. The aid reports \(E(\Q)=\Z P\), with \(P\) of infinite order, and that
\[
 \Z_{>0}\longrightarrow S,\qquad n\longmapsto x(nP)
\]
is bijective. The point at infinity is excluded by using affine \(x\)-coordinates, and the two signs of \(n\) have the same \(x\)-coordinate. In the canonical-height convention
\(\log H(x(Q))=2\widehat h(Q)+O(1)\), it gives
\[
 N_S(B)=\sqrt{\frac{\log B}{2\widehat h(P)}}+O(1)
 \qquad(B\to\infty).
\]
Thus this infinite Diophantine set has counting function smaller than every positive power of \(B\). The aid also refers to fixed-power sets as calibrations; such examples do not justify a universal positive-power lower growth threshold. The elliptic example is a concrete obstruction to that kind of interpretation of ``fast'', not a claim of an absolute thinnest set.

The conditional literature discussion has two separate hypotheses. Under Koenigsmann's stated strong Bombieri--Lang hypothesis, the aid says no infinite subset of \(\Z\) is Diophantine over \(\Q\). The exact strengthened hypothesis is not reproduced and must not be replaced by an unqualified standard Bombieri--Lang statement. Under Pasten's later approximation conjecture, it says explicit listable non-Diophantine sets follow; neither that conjecture nor those set definitions are given. Both conclusions remain conditional. They do not convert the unconditional open question into a known example.
\subsection{Short English statement}
The record asks for listable non-Diophantine rational sets and a meaningful notion of thin Diophantine sets. The aid gives an explicit elliptic-coordinate set with only square-root-logarithmic height growth, ruling out a universal positive-power growth interpretation. Its proposed non-Diophantine examples still depend on precisely specified conjectures that are not reproduced here.
\subsection{Sources}
\begin{itemize}
\item Local UnsolvedMath registry, permanent ID 20002292, exact source code catalogue entry; canonical \texttt{aim-logic-notes.json}, \texttt{notes[67]}, tag \texttt{question}; AIM Workshop Problem Lists, Logic, Question 8 (Pasten). Created and updated 14 August 2026.
\item \emph{Problems related to ``Definability and Decidability Problems in Number Theory''}, AIM workshop 9--13 September 2013, moderated by T. Scanlon with notes by J. Demeyer, according to supplied verification; \url{https://aimath.org/pastworkshops/definabilityinntproblems.pdf}. The PDF was not accessed.
\item The aid names Koenigsmann's strong Bombieri--Lang hypothesis and Pasten's later approximation conjecture, without full statements, titles, publication dates, theorem numbers, or separate URLs. It gives no complete reference for the elliptic rational-group and height calculation.
\item Attempt 1: \texttt{partial\_result}, problem status at run \texttt{open}, rigor label \texttt{proved}, novelty confidence \texttt{low}; classification \texttt{PARTIAL-PROGRESS}, registry status \texttt{partially\_solved}. Candidate: explicit height-sparse elliptic-coordinate family. The open-status and literature assessments are dated supplied evidence.
\item Ulam AI / UnsolvedMath Contributors; CC BY 4.0. Machine-generated research aid, not peer reviewed; independently verify before citation.
\end{itemize}
\end{document}
